\section{Datasets}
\done\todo{Table 1 in Chen et al.~\cite{chen2024improvingdatacurationsoftware} list a few data sets. Use this to discuss the nature and the availability of existing data sets and discuss the shortcomings for educational purpose. Discuss NVD and CVE as data sources}

Chen et al. (2024) \cite{chen2024improvingdatacurationsoftware} surveys many different publicly released collections of open source patches for various vulnerabilities. These data sets fall short in three primary attributes: language support, scale, and overall completeness.  Firstly, many datasets lack diversity in language support, and reflect a skew towards one or two languages. According to Chen et al., of the 12 datasets surveyed, 6 of them are in C/C++. Furthermore, many of these datasets contain a small project to patches ratio, with some notable cases being VulnCatcher and Big-Vul (as referenced in Chen et al.) having numbers of 3 projects to 2879 patches and 348 projects to 10547 patches, respectively. Lastly, some datasets reflect a gap in their metadata — such as VulnCode-DB, omits its project count entirely.

The shortcomings presented in the datasets surveyed in Chen et al. affect their application in the current computer science curriculum and reflect the need of a truly consolidated and organized set of curated examples. Firstly, a lack of language support across the collections of vulnerability and patch examples means that there is a large skew toward applications in one language, while the other languages lack examples to reference. Second, the small project to patches ratio for many modern datasets does not allow students to see vulnerabilities in different contexts, beyond the scope of a few projects and their nuances. Lastly, the gap in some datasets' metadata means some datasets lack critical information that would otherwise enable a user to gauge diversity. This greatly reduces the organization that would be useful for educational purposes as you lack the information to organize patches by projects or are unable to distinguish patches due to the lack of a total project count.

The datasets considered in Table I of Chen et al. derive data from the NVD and CVE as sources. The CVE provides identifiers and descriptions of every known vulnerability. The NVD uses CVE entries and provides descriptive details — most importantly, a section for links to patches for vulnerabilities. Because the NVD sources patches from open-source projects, they are susceptible to issues. NVD-sourced data is therefore rife with issues such as broken links and varying quality of patches. For example, of 6,628 CVEs, 66.57\% contained references to patches, of which 67.21\% were functional. The issues of the datasets' data sources propagate into the datasets themselves — numerous projects and fixes are excluded leading to the shortcomings seen in Table I of Chen et al., further emphasizing the need for a more robust and curated collection of examples.